\subsection{Sensordaten in IoT-Systemen}
Im Kontext von \gls{iot}-Anwendungen werden Sensordaten kontinuierlich erfasst und an zentrale Systeme wie Datenbanken oder Cloud-Plattformen übertragen. 
Diese Daten bilden die Grundlage für weiterführende Analysen und die Ableitung anwendungsspezifischer Funktionen.
\gls{iot}-Systeme erfassen dabei eine Vielzahl physikalischer und umgebungsbezogener Größen, darunter Temperatur, Standortdaten und weitere Umweltparameter. 
Diese Informationen werden durch Sensoren bereitgestellt und bilden die Basis für kontextabhängige Anwendungen~\cite{7123563}. 

Die erfassten Sensordaten werden analysiert und zur Ableitung automatisierter Entscheidungen verwendet. 
Dadurch lassen sich intelligente Dienste realisieren, die auf aktuellen Zustandsinformationen basieren. 
Neben lokal erfassten Sensordaten können auch externe Datenquellen in \gls{iot}-Systeme integriert werden, um zusätzliche Informationen bereitzustellen. 
Beispielsweise lassen sich in Smart-Home-Anwendungen Aktionen auf Basis von Wetterinformationen auslösen~\cite{7123563}. 

Typischerweise erfolgt die Erfassung der Sensordaten in festen Zeitintervallen. 
Die gewonnenen Messwerte werden anschließend im System weiterverarbeitet und für die jeweilige Anwendung aufbereitet.